% !TEX root = ../main.tex
\chapter{Representation theory of Lie groups and Lie algebras}\label{ch:l17}
\begin{paracol}{2}

\begin{sources}
\begin{tabularx}{\linewidth}{@{}l l L@{}}
\toprule
\textbf{Video} & \textbf{Length} & \textbf{Content}\\
\midrule
\seg{17}{1}{Lecture 17: Representation theory of Lie groups and Lie algebras} & 1:32:39 & A short section added on request (``mostly you meet Lie algebras and groups in physics in terms of their representations''), without proofs: linear representations of a Lie algebra and the representation space; examples ($\mathfrak{sl}(2,\C)$ on $\C^2$, the rotation algebra on $\R^3$ and on $\C^2$, the trivial and the adjoint representation); reducible, irreducible and faithful representations; the Killing form of a representation, the Casimir operator, Schur's lemma and the value $\dim L/\dim V$, checked on the vector and spin representations (with a long digression on a factor $2$); representations of a Lie group, the example $SO(2)$ and the adjoint representation of a Lie group.\\
\bottomrule
\end{tabularx}

\smallskip
\textbf{Transcript} \file{materials/transcripts/C17-V01.txt} (YouTube automatic captions; the blackboard matrices are reconstructed from the verbal description, several signs are unsure already in the lecture and are checked in the course notes).
\end{sources}

\begin{reading}[Where this lecture sits in the other notes]
\emph{Quantum Theory}, Lectures 13--15; \emph{Theoretische Quantenmechanik}, Lectures 14--16 (angular momentum, Clebsch--Gordan): forward pointers, since the spin representations of $\mathfrak{so}(3)$ and the addition of angular momenta are the physics of the examples below [check later]. It uses the algebras of \cref{ch:l13,ch:l14,ch:l16}; the next lecture goes back from the algebra to the group (\cref{ch:l18}).
\end{reading}

% ==================================================================
\section{Representations of a Lie algebra}\label{sec:rp-algebra}

After the previous lecture some listeners asked about representations of Lie algebras and groups, and the lecturer decided to add a brief section, because in physics one mostly meets Lie algebras and Lie groups \emph{through} their representations, and very often they are even defined by them \ts{17}{1}{0:00} \ts{17}{1}{0:21}. In this course a Lie group was defined as a smooth manifold with a group law compatible with the smooth structure, and its Lie algebra as the left-invariant vector fields (which one can pull back to $T_{\mathbb 1}G$ with an induced bracket, \cref{thm:ll-iso}). ``But what on earth are these representations?'' \ts{17}{1}{0:41}. A half-semester course could be given on the subject, so the lecture presents \emph{no proofs}, only the important concepts \ts{17}{1}{1:02} \ts{17}{1}{1:28} \ts{17}{1}{1:49}.

\begin{definition}[New concepts: representation of a Lie algebra]\label{def:rp-rep}
Let $(L,[\cdot,\cdot])$ be a Lie algebra. A \textbf{representation} of $L$ is a linear map
\[
  \rho\colon L\to\mathrm{End}(V)
\]
into the endomorphisms of some vector space $V$ over $\R$ or $\C$ (``linear'': $L$ is a vector space, and $\mathrm{End}(V)$ is a vector space), such that for all $a,b\in L$
\[
  \rho\bigl([a,b]\bigr)=\bigl[\rho(a),\rho(b)\bigr]_{\mathrm{End}(V)} .
\]
The vector space $V$ is the \textbf{representation space}; if $V$ is finite-dimensional the representation is finite-dimensional. Since the maps $\rho(a)$ are linear maps on $V$ (not vector fields), the bracket on the right is not the differential-geometric one but the commutator, $[\phi,\psi]:=\phi\circ\psi-\psi\circ\phi$. This is bilinear, antisymmetric by construction and satisfies the Jacobi identity: this is why one maps into $\mathrm{End}(V)$, which thanks to its natural composition is a Lie algebra at once. So: one may first take the bracket in $L$ and then map, or map first and then take the commutator, with the same result \ts{17}{1}{2:10} \ts{17}{1}{2:52} \ts{17}{1}{3:34} \ts{17}{1}{4:17} \ts{17}{1}{4:59} \ts{17}{1}{5:43}.
\end{definition}
The terminology can confuse: one might think that $\mathrm{End}(V)$ is the ``space where $L$ is represented'', but the name \emph{representation space} is reserved for $V$. The idea is that, once $\rho(a)$ is known, it can \emph{act on} the vectors of $V$: an endomorphism takes a vector and produces a new vector \ts{17}{1}{6:03} \ts{17}{1}{6:25} \ts{17}{1}{7:07}. (``Linear'' in ``linear representation'' refers to the fact that the elements are mapped to endomorphisms of a vector space.) The name comes from this: the abstract elements of $L$ are \emph{represented} by concrete linear maps (matrices) \ts{17}{1}{7:28}.
\begin{coursenote}[Not said in class: real algebra, complex space]
$L$ and $V$ need not be over the same field: a real Lie algebra such as $\mathfrak{so}(3,\R)$ may be represented on a complex space such as $\C^2$ (then $\rho$ is $\R$-linear), as is done below. \ts{17}{1}{19:37}
\end{coursenote}

% ------------------------------------------------------------------
\section{Examples}\label{sec:rp-examples}

\subsection{$\mathfrak{sl}(2,\C)$ on $\C^2$}\label{sec:rp-sl2}
The algebra of \cref{ch:l15,ch:l16} has three basis vectors with
\[
  [X_1,X_2]=2X_2,\qquad[X_1,X_3]=-2X_3,\qquad[X_2,X_3]=X_1 ,
\]
derived from the left-invariant vector fields; so far $X_1,X_2,X_3$ were only the representatives at the identity of those fields, ``we didn't know more'' \ts{17}{1}{7:49} \ts{17}{1}{8:31}.
\begin{recall}[The three brackets]
The relations are those of \cref{thm:sa-bracket}.
\end{recall}
A representation $\rho\colon\mathfrak{sl}(2,\C)\to\mathrm{End}(\C^2)$ is obtained by prescribing it on a basis (linearity does the rest) \ts{17}{1}{8:52} \ts{17}{1}{9:14}:
\[
  \rho(X_1)=\begin{pmatrix}1&0\\0&-1\end{pmatrix},\quad
  \rho(X_2)=\begin{pmatrix}0&1\\0&0\end{pmatrix},\quad
  \rho(X_3)=\begin{pmatrix}0&0\\1&0\end{pmatrix}.
\]
Here $\mathrm{End}(\C^2)$ is written by matrices with respect to the canonical basis, and the composition of endomorphisms is matrix multiplication. The compatibility condition has to be checked on the pairs of basis vectors, because both sides are bilinear in $a,b$. For $(X_1,X_2)$:
\[
  \rho(X_1)\rho(X_2)-\rho(X_2)\rho(X_1)=\begin{pmatrix}0&1\\0&0\end{pmatrix}-\begin{pmatrix}0&-1\\0&0\end{pmatrix}=2\begin{pmatrix}0&1\\0&0\end{pmatrix}=2\rho(X_2)=\rho(2X_2)=\rho([X_1,X_2]),
\]
by linearity of $\rho$ \ts{17}{1}{10:15} \ts{17}{1}{10:35} \ts{17}{1}{11:18} \ts{17}{1}{12:00}. The pairs $(X_1,X_3)$ and $(X_2,X_3)$ are checked in the same way (the lecture does not carry this out; one finds $-2\rho(X_3)$ and $\rho(X_1)$) \ts{17}{1}{12:21}. So this is a representation, and it is how $\mathfrak{sl}(2,\C)$ is usually presented.

The image $\rho(\mathfrak{sl}(2,\C))\subset\mathrm{End}(\C^2)$ is the span of the three matrices, that is, the matrices $\begin{pmatrix}\alpha&\beta\\\gamma&\delta\end{pmatrix}$ with $\alpha+\delta=0$: the \emph{trace-free} matrices. This is the origin of the usual sentence ``the group $SL(2,\C)$ consists of the matrices of determinant $1$, and its algebra of the trace-free $2\times2$ matrices'': it is just a statement about \emph{one} representation. Here one went the abstract way, from the group to the left-invariant fields, and the algebra is not a priori a space of matrices. And the point of the whole business is that a given Lie algebra has, under suitable conditions, \emph{several} representations; this is only one of them \ts{17}{1}{12:42} \ts{17}{1}{13:45} \ts{17}{1}{14:05}.

\subsection{The rotation algebra on $\R^3$ and on $\C^2$}\label{sec:rp-so3}
The algebra of the rotation group appears in physics all the time \ts{17}{1}{14:27}. The algebra $\mathfrak{so}(3,\R)$ is three-dimensional and is defined, like any Lie algebra, by the choice of a basis $X_1,X_2,X_3$ and of the structure constants (\cref{def:cl-structure}). The lecture does it as follows: first the Killing form $K_{ij}=C^m{}_{in}C^n{}_{jm}$ (\cref{def:cl-killing}); then the first index of the structure constants is pulled down with it, $C_{kij}:=K_{km}C^m{}_{ij}$ (``one could say it basis-free, but that is just a way of saying this''); and the algebra $\mathfrak{so}(3,\R)$ is the one for which
\[
  C_{kij}=\varepsilon_{ijk}\in\{+1,-1,0\}
\]
(even permutation of $123$, odd permutation, no permutation: for instance $112\mapsto0$) \ts{17}{1}{14:48} \ts{17}{1}{15:08} \ts{17}{1}{15:50} \ts{17}{1}{16:31} \ts{17}{1}{16:52}. These are totally antisymmetric structure coefficients, and to be able to say \emph{totally} antisymmetric one needs the Killing form to pull the index down; conversely, from $\varepsilon_{ijk}$ one reconstructs the $C^k{}_{ij}$ \ts{17}{1}{17:12}.
\begin{coursenote}[Not said in class: the normalisation]
Taken literally, the two conditions are not compatible with unit constants: for $C^k{}_{ij}=c\,\varepsilon_{ijk}$ one has $K_{ij}=-2c^2\delta_{ij}$, hence $K_{km}C^m{}_{ij}=-2c^3\varepsilon_{ijk}$. What is meant, and what is used afterwards, is the standard algebra $[X_i,X_j]=\varepsilon_{ijk}X_k$, whose Killing form is $K=-2\delta$ (negative definite). \ts{17}{1}{16:52}
\end{coursenote}
\textbf{On $\R^3$} (the vector representation) one takes $\rho_V\colon\mathfrak{so}(3,\R)\to\mathrm{End}(\R^3)$ with three antisymmetric matrices, each with a pair of $\pm1$ entries and zeros elsewhere \ts{17}{1}{17:33} \ts{17}{1}{17:53} \ts{17}{1}{18:14}. The lecture does not trust its signs (``up to signs this is right; please check for yourself'') and says one has to check the compatibility condition; the consistent choice is
\[
  \bigl(\rho_V(X_i)\bigr)_{jk}=-\varepsilon_{ijk},\quad\text{e.g. }\rho_V(X_1)=\begin{pmatrix}0&0&0\\0&0&-1\\0&1&0\end{pmatrix}
\]
(\emph{not said in class}: with this choice $[\rho_V(X_i),\rho_V(X_j)]=\varepsilon_{ijk}\rho_V(X_k)$.) \ts{17}{1}{18:35} \ts{17}{1}{18:55}. This is what one learns first: the rotation algebra acts on $\R^3$.

\textbf{On $\C^2$} (the spin representation): famously from quantum mechanics there is another one, $\rho_{\mathrm{spin}}\colon\mathfrak{so}(3,\R)\to\mathrm{End}(\C^2)$, which acts on two-component complex objects, the spinors of Schr\"odinger theory, and the matrices are $\frac12$ of the Pauli matrices $\sigma_1,\sigma_2,\sigma_3$ ``with or without a factor $\mathrm{i}$; please figure it out'' \ts{17}{1}{19:17} \ts{17}{1}{19:57} \ts{17}{1}{20:18}. The consistent choice is worked out later in the lecture (\cref{sec:rp-casimir-spin}): $\rho_{\mathrm{spin}}(X_k)=-\frac{\mathrm i}{2}\sigma_k$. So the \emph{same} rotation algebra can be represented on $\R^3$ (rotating in real three-space) or on $\C^2$ (spinors), and these are different representations \ts{17}{1}{20:39} \ts{17}{1}{21:00}.

\subsection{The trivial and the adjoint representation}\label{sec:rp-triv}
\textbf{Trivial representation} \ts{17}{1}{21:21} \ts{17}{1}{21:41}: for every Lie algebra and on every vector space $V$ there is $\rho_{\mathrm{triv}}\colon L\to\mathrm{End}(V)$, $a\mapsto0$ (the zero map, ``the matrix with only zeros''). It is a representation: $\rho([a,b])=0$ and $[\rho(a),\rho(b)]=[0,0]=0$, the same zero of $\mathrm{End}(V)$, so compatibility always holds \ts{17}{1}{22:09} \ts{17}{1}{22:29}. Physicists call it the \emph{scalar} representation (nothing changes under the algebra) \ts{17}{1}{23:11}.

\textbf{Adjoint representation} \ts{17}{1}{23:31}. Does every Lie algebra have a \emph{non-trivial} representation? Yes. A representation has two choices: the representation space, and, once the space is fixed, still the linear map, so different representations may live on the same space \ts{17}{1}{23:52} \ts{17}{1}{24:13} \ts{17}{1}{24:54}. If one is short of vector spaces, one is already given one: the vector space underlying $L$. Take $V=L$ and
\[
  \mathrm{ad}\colon L\to\mathrm{End}(L),\qquad a\mapsto\mathrm{ad}_a=[a,\,\cdot\,],
\]
that is, $\mathrm{ad}_a(\ell)=[a,\ell]\in L$ (\cref{def:cl-ad}). Compatibility is the Jacobi identity. This is the \textbf{adjoint representation} \ts{17}{1}{25:14} \ts{17}{1}{26:05} \ts{17}{1}{26:27} \ts{17}{1}{26:49}. The whole classification of semisimple and simple Lie algebras (\cref{ch:l14}) made essential use of it: the Cartan subalgebra is where $\mathrm{ad}$ is diagonalised, and all the other directions are eigenvectors \ts{17}{1}{27:10}.

% ------------------------------------------------------------------
\section{Reducible, irreducible and faithful representations}\label{sec:rp-irred}

To study representations one wants \emph{building blocks}: if a representation can be understood through two others, there is no need to study it separately \ts{17}{1}{27:30}.
\begin{definition}[New concepts: reducible and irreducible representations]\label{def:rp-irred}
A representation $\rho\colon L\to\mathrm{End}(V)$ is \textbf{reducible} if there exists a vector subspace $U\subseteq V$, with $U\neq\{0\}$ and $U\neq V$, such that
\[
  \rho(a)u\in U\qquad\text{for all }a\in L,\ u\in U .
\]
Such a $U$ is an \textbf{invariant subspace}. Equivalently $\rho$ restricts to a representation $\rho_U\colon L\to\mathrm{End}(U)$. Otherwise $\rho$ is called \textbf{irreducible}. (``Reducible'': it can be reduced to the smaller pieces $U$; ``irreducible'': it cannot.) \ts{17}{1}{27:51} \ts{17}{1}{28:12} \ts{17}{1}{28:53} \ts{17}{1}{29:13} \ts{17}{1}{29:54}
\end{definition}
\begin{answer}{Should the invariant subspace be non-trivial?}[Proper invariant subspace]
A student asks about the condition ``$U$ non-zero'': it does not suffice to exclude $\{0\}$, which is always invariant; the whole space $V$ must be excluded too. The lecturer agrees (``I should have written non-trivial''), and the definition above contains both conditions. \ts{17}{1}{36:52}
\end{answer}
\textbf{Example} \ts{17}{1}{30:15}. Take the two representations $\rho_V\colon\mathfrak{so}(3,\R)\to\mathrm{End}(\R^3)$ and $\rho_{\mathrm{spin}}\colon\mathfrak{so}(3,\R)\to\mathrm{End}(\C^2)$ and build a new one on the direct sum, $V=\R^3\oplus\C^2$ (the real $\R^3$ is considered as inside $\C^3$, so the whole is a complex matrix representation): to $a$ one associates the block matrix
\[
  \rho(a)=\begin{pmatrix}\rho_V(a)&0\\0&\rho_{\mathrm{spin}}(a)\end{pmatrix}
\]
with zeros off the diagonal blocks \ts{17}{1}{30:56} \ts{17}{1}{31:17}. Matrix multiplication can be carried out blockwise, so the commutator of two such matrices has the same block form with the commutators of the blocks; since each block is a representation, $\rho$ is again one \ts{17}{1}{32:20} \ts{17}{1}{32:41}. Now $U=\R^3$ (the vectors with the last two entries zero) and $U'=\C^2$ (the first three entries zero) are vector subspaces and \emph{invariant}: $\rho(a)$ acts on a vector with entries only in the upper block and, because of the block form, nothing is produced in the lower one \ts{17}{1}{33:04} \ts{17}{1}{33:46}. So a representation with a block form is reducible. One can also show (not done in class) that $\rho_V$ and $\rho_{\mathrm{spin}}$ themselves are not reducible \ts{17}{1}{34:07}.
\begin{coursenote}[Not said in class: building blocks]
``One can always understand a reducible representation through irreducible building blocks'' holds for finite-dimensional representations of \emph{semisimple} Lie algebras over $\C$ (complete reducibility, Weyl's theorem): every such representation is a direct sum of irreducible ones. For a general Lie algebra a reducible representation need not be a direct sum, only an invariant subspace with a quotient (a block-\emph{triangular} form). \ts{17}{1}{34:28}
\end{coursenote}
\begin{aside}[Elementary particles]
Physics: ``particles are irreducible representations of the symmetry'' \ts{17}{1}{34:48}. For the particle of classical mechanics the symmetry is the Galilei group (rotations, Galilei boosts, translations in space-time); to keep it simple, one traps the particle at a point: what remains is rotational invariance, with the group $SO(3)$ and the algebra above. The particle may then be represented by a vector, by a scalar (the trivial representation) or, in quantum mechanics, by a two-component spinor. If one took a reducible representation one could always make \emph{two} particles out of it (two particles also transform in some representation of the symmetry, and one wants to study each on its own); an \emph{elementary} particle is an irreducible one. This is just the reason why the term comes up in particle physics \ts{17}{1}{35:08} \ts{17}{1}{35:50} \ts{17}{1}{36:10}.
\end{aside}
\begin{definition}[New concepts: faithful representation]\label{def:rp-faithful}
A representation $\rho\colon L\to\mathrm{End}(V)$ is \textbf{faithful} if $\rho$ is injective: no two different elements of $L$ are mapped to the same endomorphism. (The representation then ``represents $L$ without loss'': $L$ can be read off from its image.) \ts{17}{1}{37:13} \ts{17}{1}{37:35} \ts{17}{1}{38:33}
\end{definition}
\begin{exercise}[Adjoint representation: irreducible, faithful?]\label{ex:rp-ad}
Check irreducibility and faithfulness of the adjoint representation $\rho=\mathrm{ad}$ (the lecture: ``let us formulate this as a homework problem''). \emph{Hint, not said in class:} an invariant subspace of $\mathrm{ad}$ is exactly an ideal, and $\ker\mathrm{ad}$ is the centre of $L$; see \cref{def:cl-simple,thm:sa-simple}. (\cref{exam:ch17}.) \ts{17}{1}{38:54}
\end{exercise}

% ------------------------------------------------------------------
\section{The Casimir operator}\label{sec:rp-casimir}

A very important concept in the study of representations is the \emph{Casimir operator} (the name is not explained in class; not said in class: it is named after the Dutch physicist Hendrik Casimir) \ts{17}{1}{39:40}. In physics one is very often handed the Lie algebra already in a representation, so there is some confusion: a Casimir operator cannot be constructed for a Lie algebra, but for \textbf{any given representation of any given Lie algebra} \ts{17}{1}{40:00} \ts{17}{1}{40:23}. This takes several steps.

\textbf{Step A: the Killing form of a representation.} \ts{17}{1}{40:44} The space $\mathrm{End}(V)$ with the commutator is a Lie algebra, so it has a Killing form, but the lecture wants something different, a form on $L$ itself \ts{17}{1}{41:05} \ts{17}{1}{41:47}.
\begin{definition}[New concepts: the Killing form of a representation]\label{def:rp-killing}
For a representation $\rho\colon L\to\mathrm{End}(V)$ the \textbf{Killing form with respect to $\rho$} is
\[
  K_\rho\colon L\times L\to\C\ (\text{or }\R),\qquad K_\rho(a,b):=\mathrm{tr}_V\bigl(\rho(a)\circ\rho(b)\bigr).
\]
For $\rho=\mathrm{ad}$ this is the Killing form $K(a,b)=\mathrm{tr}(\mathrm{ad}_a\circ\mathrm{ad}_b)$ used so far (\cref{def:cl-killing}): that one was the Killing form of the \emph{adjoint} representation. \ts{17}{1}{42:07} \ts{17}{1}{42:50} \ts{17}{1}{43:12}
\end{definition}
\textbf{Step B: the Casimir operator} \ts{17}{1}{43:36} \ts{17}{1}{43:57}.
\begin{definition}[New concepts: the Casimir operator]\label{def:rp-casimir}
Let $\rho\colon L\to\mathrm{End}(V)$ be a representation. Choose a basis $X_1,\dots,X_{\dim L}$ of $L$ and let $\tilde X_1,\dots,\tilde X_{\dim L}$ be the \emph{dual basis} of $L$ with respect to $K_\rho$, that is
\[
  K_\rho\bigl(X_i,\tilde X_j\bigr)=\delta_{ij}.
\]
The \textbf{Casimir operator} of $\rho$ is
\[
  \Omega_\rho:=\sum_{i=1}^{\dim L}\rho(X_i)\circ\rho(\tilde X_i)\ \in\mathrm{End}(V).
\]
It is independent of the choice of the basis $X_i$. \ts{17}{1}{44:18} \ts{17}{1}{45:02} \ts{17}{1}{45:24} \ts{17}{1}{45:44} \ts{17}{1}{46:05}
\end{definition}
\begin{coursenote}[Not said in class: existence of the dual basis, signs]
The dual basis exists iff $K_\rho$ is non-degenerate. The lecture first writes $\tilde X_i$ ``for an orthonormal basis'', $\tilde X_i=X_i$, and says this works only if the group is compact (``otherwise I do not find such a basis''). For a compact group $K$ is \emph{negative} definite, so an orthonormal basis has $K_\rho(X_i,X_j)=-\delta_{ij}$ and one writes $\Omega_\rho=-\sum_i\rho(X_i)\rho(X_i)$; the lecture corrects this sign later \ts{17}{1}{1:00:08} \ts{17}{1}{1:00:28}. With the dual basis of \cref{def:rp-casimir} no compactness is needed, and the signs take care of themselves (see the examples). Basis independence: under $X_i'=M_i{}^jX_j$ the dual basis transforms with $(M^{-1})^T$ and the sum is invariant. \ts{17}{1}{46:25} \ts{17}{1}{46:50}
\end{coursenote}
The Casimir operator can be applied to a vector of the representation space and gives a vector of the same space \ts{17}{1}{47:11}. The first property:
\begin{theorem}[The Casimir operator commutes with the representation]\label{thm:rp-commute}
For every $a\in L$: $[\Omega_\rho,\rho(a)]=0$. Informally: it commutes with everything in the image $\rho(L)$.
\end{theorem}
The lecture does not prove it and asks for the proof on the problem sheet (\cref{exam:ch17}) \ts{17}{1}{47:32} \ts{17}{1}{47:52} \ts{17}{1}{48:13}. This is how the operator is usually introduced \ts{17}{1}{48:36}. The second statement is \textbf{Schur's lemma} (the transcript has ``shortlemma''; named after Issai Schur, not said in class) \ts{17}{1}{49:00}:
\begin{theorem}[Schur's lemma for the Casimir operator]\label{thm:rp-schur}
If $\rho$ is irreducible, then there is a number $c_\rho\in\C$ (or $\R$) such that
\[
  \Omega_\rho=c_\rho\,\mathrm{id}_V ,
\]
a diagonal matrix with $c_\rho$ on the diagonal. (The lecture gives no proof.) \ts{17}{1}{49:21}\ts{17}{1}{50:03}
\end{theorem}
So the Casimir operator of an irreducible representation is a constant multiple of the identity. Not said in class: over $\R$ the lemma can fail (an irreducible real representation may have commutant larger than the multiples of the identity, e.g.\ the rotation $\begin{pmatrix}0&-1\\1&0\end{pmatrix}$ commutes with the irreducible $SO(2)$ action on $\R^2$); the proof needs $\C$ (an eigenvalue exists, and its eigenspace is invariant).

The lecture then tries to show that this number is always given by the ratio of the dimensions \ts{17}{1}{50:25} \ts{17}{1}{50:45}:
\begin{theorem}[The value of the Casimir operator]\label{thm:rp-value}
If $\rho\colon L\to\mathrm{End}(V)$ is irreducible, $\dim V\ne0$, and $K_\rho$ is non-degenerate, then
\[
  c_\rho=\frac{\dim L}{\dim V}.
\]
\end{theorem}
\begin{proof}
By \cref{thm:rp-schur}, $\Omega_\rho=c_\rho\,\mathrm{id}_V$, so taking the trace (linear, and $\mathrm{tr}\,\mathrm{id}_V=\dim V$):
\[
  \mathrm{tr}\,\Omega_\rho=c_\rho\,\dim V .
\]
On the other hand, by linearity of the trace and the definitions,
\[
  \begin{aligned}
  \mathrm{tr}\,\Omega_\rho&=\sum_i\mathrm{tr}\bigl(\rho(X_i)\rho(\tilde X_i)\bigr)\\
  &=\sum_iK_\rho(X_i,\tilde X_i)\\
  &=\sum_{i=1}^{\dim L}1=\dim L .
  \end{aligned}
\]
Hence $c_\rho\dim V=\dim L$. \ts{17}{1}{51:15} \ts{17}{1}{51:35} \ts{17}{1}{52:16} \ts{17}{1}{52:37} \ts{17}{1}{53:01} \ts{17}{1}{53:43} \ts{17}{1}{54:04}
\end{proof}
So different irreducible representations of the same Lie algebra on different representation spaces $V_1,V_2$ are labelled by this number, and one can act with the Casimir operator on any vector of the representation space of one particular irreducible representation and always get the same number times the vector \ts{17}{1}{54:29} \ts{17}{1}{54:50}.

\subsection{Check: the vector representation}\label{sec:rp-casimir-vec}
In physics the rotation algebra reads $[J_i,J_j]=\varepsilon_{ijk}J_k$ (``something of that type''; here the abstract $X_i$) and $\vec J^2$ is not literally a square of a vector but $\sum_i\rho(J_i)\rho(J_i)$ \ts{17}{1}{55:10} \ts{17}{1}{55:31} \ts{17}{1}{56:12}. Take the vector representation. Each $\rho_V(X_i)^2$ is a diagonal matrix with $-1$ at the two positions in its block, so the three squares add up to $-2\cdot\mathbb 1_3$ (the lecture phrases it with the rotation blocks of the $\pm1$ entries and gets the pattern $(1,1,0)$, $(1,0,1)$, $(0,1,1)$ for the three squares, so that the sum is $2\cdot\mathbb 1$) \ts{17}{1}{56:33} \ts{17}{1}{57:16} \ts{17}{1}{57:37} \ts{17}{1}{57:57}. With the prefactor $\frac12$ this gives $c_{\rho_V}=1$, and \cref{thm:rp-value} agrees: $\dim\mathfrak{so}(3)/\dim\R^3=3/3=1$ \ts{17}{1}{58:18} \ts{17}{1}{58:39} \ts{17}{1}{59:00} \ts{17}{1}{59:25}. The lecture stresses that in physics the factor $\frac12$ is absent, and everything is then off by a factor $2$.
\begin{answer}{Where is the sign?}[The minus sign in the Casimir operator]
A student remarks on a sign in the computation (the question is inaudible in the transcript [?]; it concerns the minus signs coming from the squares of the blocks). The lecturer answers that this is the earlier problem with the sign: for a compact group the Killing form is negative definite and the definition of $\Omega_\rho$ needs the overall minus sign, ``then the signs are right''. See the course note after \cref{def:rp-casimir}. \ts{17}{1}{59:48} \ts{17}{1}{1:00:08} \ts{17}{1}{1:00:48}
\end{answer}

\subsection{Check: the spin representation, and a factor $2$}\label{sec:rp-casimir-spin}
Now the second example, the same abstract algebra $[X_i,X_j]=\varepsilon_{ijk}X_k$ in the spin representation $\rho_{\mathrm{spin}}$ \ts{17}{1}{1:01:09}. The Pauli matrices are (the lecturer needs reminding)
\[
  \sigma_1=\begin{pmatrix}0&1\\1&0\end{pmatrix},\quad\sigma_2=\begin{pmatrix}0&-\mathrm i\\\mathrm i&0\end{pmatrix},\quad\sigma_3=\begin{pmatrix}1&0\\0&-1\end{pmatrix}.
\]
Check of $[\sigma_1,\sigma_2]$: $\sigma_1\sigma_2=\mathrm i\sigma_3$ and $\sigma_2\sigma_1=-\mathrm i\sigma_3$, so $[\sigma_1,\sigma_2]=2\mathrm i\sigma_3$ \ts{17}{1}{1:01:30} \ts{17}{1}{1:01:50} \ts{17}{1}{1:02:14} \ts{17}{1}{1:02:37}. The factor $2$ is removed by taking $\tau_i=\frac12\sigma_i$, for which $[\tau_1,\tau_2]=\mathrm i\tau_3$; the factor $\mathrm i$ by putting $-\mathrm i$ in front, since $(-\mathrm i)^2[\tau_1,\tau_2]=-\mathrm i\tau_3$ is exactly the representative of $X_3$ (the lecturer counts the minus signs aloud and arrives at this choice). In summary \ts{17}{1}{1:02:57} \ts{17}{1}{1:03:19} \ts{17}{1}{1:03:40} \ts{17}{1}{1:04:02} \ts{17}{1}{1:04:22} \ts{17}{1}{1:04:42}:
\[
  \rho_{\mathrm{spin}}(X_k)=-\frac{\mathrm i}{2}\sigma_k,\qquad[\rho_{\mathrm{spin}}(X_i),\rho_{\mathrm{spin}}(X_j)]=\varepsilon_{ijk}\rho_{\mathrm{spin}}(X_k).
\]
(In quantum mechanics one wants hermitian generators and puts an extra $\mathrm i$ in.) The Casimir operator is then computed with the same recipe, $\Omega=-\frac12\sum_i\rho(X_i)\rho(X_i)$ \ts{17}{1}{1:05:02} \ts{17}{1}{1:05:23}. The lecture first evaluates it with the matrices $-\mathrm i\sigma_i$ (the factor $\frac12$ inside is forgotten): each square is $-\mathbb 1$, ``the minuses all go out and there is a three everywhere'', so $\Omega=\frac32\mathbb 1$, in agreement with $\dim L/\dim V=\frac32$ (the lecture also writes it as $\frac12(\frac12+1)$, which is $\frac34$: a slip, see the course note below) \ts{17}{1}{1:05:44} \ts{17}{1}{1:06:26} \ts{17}{1}{1:07:07} \ts{17}{1}{1:07:30} \ts{17}{1}{1:07:50}.

Then something is wrong: with the correct matrices $-\frac{\mathrm i}2\sigma_i$ each square is $-\frac14\mathbb 1$ and one gets $\frac12\cdot3\cdot\frac14=\frac38$, not $\frac32$ \ts{17}{1}{1:08:12} \ts{17}{1}{1:08:32} \ts{17}{1}{1:08:55}. And $\frac38$ is not $j(j+1)$ either (``in the vector case it wasn't either'': for the vector representation one has $j=1$, $j(j+1)=2$ but the number found is $1$). So the lecture's first diagnosis is a mismatch of conventions between physics and mathematics, by an overall factor $2$: the physical $\vec J^2$ is \emph{twice} the Casimir operator defined here, ``because one wants the dual basis with respect to the Killing form on the same space, to make everything basis independent'' \ts{17}{1}{1:09:16} \ts{17}{1}{1:09:36} \ts{17}{1}{1:09:57} \ts{17}{1}{1:10:18} \ts{17}{1}{1:10:38} \ts{17}{1}{1:10:59} \ts{17}{1}{1:11:20}.
\begin{answer}{Why does the number not come out as $3/2$?}[A student: the basis depends on $\rho$]
After a pause for thinking (``I do not get the algebra straight'') a student points out the origin of the trouble: the dual vectors $\tilde X_i$ were defined by $K_\rho(X_i,\tilde X_j)=\delta_{ij}$ (the lecture: by $-\delta_{ij}$ on a compact group), and $K_\rho(X_i,X_j)=\mathrm{tr}(\rho(X_i)\rho(X_j))$ \emph{depends on the representation}. In the spin representation the proportionality constant between $\tilde X_i$ and $X_i$ is therefore different from the one in the vector representation, and the earlier computation, done with the constant of the vector representation, does not apply. The lecturer thanks him (``it's an excellent problem-sheet question'': the dual basis is constructed with respect to $K_\rho$, so it has to be adapted to every $\rho$), and says that it is now clear that the missing factor $2$ should be there, so that one gets $\frac32$, but ``I should make no further guesses, just calculate''. \ts{17}{1}{1:11:50} \ts{17}{1}{1:13:46} \ts{17}{1}{1:14:07} \ts{17}{1}{1:14:50} \ts{17}{1}{1:15:32} \ts{17}{1}{1:15:52} \ts{17}{1}{1:16:13}
\end{answer}
\begin{coursenote}[Not said in class: the computation, and both factors of 2]
With $\rho_{\mathrm{spin}}(X_k)=-\frac{\mathrm i}{2}\sigma_k$ one has $K_{\rho_{\mathrm{spin}}}(X_i,X_j)=\mathrm{tr}\bigl(-\tfrac14\sigma_i\sigma_j\bigr)=-\tfrac12\delta_{ij}$. The dual basis is therefore $\tilde X_i=-2X_i$ and
\[
  \Omega_{\mathrm{spin}}=\sum_i\rho(X_i)\rho(\tilde X_i)=-2\sum_i\Bigl(-\tfrac14\Bigr)\mathbb 1=\tfrac32\,\mathbb 1=\frac{\dim L}{\dim V}\mathbb 1 .
\]
For the vector representation $K_{\rho_V}(X_i,X_j)=-2\delta_{ij}$, $\tilde X_i=-\frac12X_i$ and $\Omega_V=\mathbb 1=\frac33\mathbb 1$ (as above). So the class has the right numbers when the dual basis is taken \emph{with respect to $K_\rho$}. If instead one uses always the \emph{same} basis, dual with respect to the adjoint Killing form $K=-2\delta$ ($\tilde X_i=-\frac12X_i$ for every $\rho$), then $-\frac12\sum_i\rho(X_i)^2=\frac38$ in the spin and $1$ in the vector representation, that is $\frac12j(j+1)$ with $j=\frac12,1$: this is the lecture's ``factor $2$ with respect to physics'', where $\vec J^2=-\sum_i\rho(X_i)\rho(X_i)=j(j+1)$ ($\frac34$ and $2$). So both statements of the lecture are right, they refer to two different normalisations. (The first evaluation, $\frac32=\frac12(\frac12+1)$, was a slip: $\frac12(\frac12+1)=\frac34$ and $\frac32=2j(j+1)$.) \ts{17}{1}{1:09:36}
\end{coursenote}
So different representation spaces, at least when they differ in dimension, are labelled by different numbers in front of the identity in the Casimir operator, and one says: the particle has spin $\frac12$ or spin $1$, and the Casimir operator gives the same value on any representing vector. The lecturer still maintains the factor $2$ between mathematics and physics, but the other comment is important \ts{17}{1}{1:16:35} \ts{17}{1}{1:16:57} \ts{17}{1}{1:17:17}.

% ------------------------------------------------------------------
\section{Representations of a Lie group}\label{sec:rp-group}

Now Lie groups \ts{17}{1}{1:17:44}.
\begin{definition}[New concepts: representation of a Lie group]\label{def:rp-grouprep}
Let $(G,\bullet)$ be a finite-dimensional Lie group. A \textbf{representation} of $G$ on a vector space $V$ is a map
\[
  R\colon G\to GL(V),
\]
where $GL(V)\subseteq\mathrm{End}(V)$ is the set of linear maps $\varphi\colon V\to V$ with $\det\varphi\ne0$, that is, the \emph{invertible} endomorphisms (over $\C^n$ or $\R^n$ this is the general linear group). Since the composition of invertible maps is invertible, $\circ$ is a group law on $GL(V)$, and one requires compatibility with the group operations:
\[
  R(g_1\bullet g_2)=R(g_1)\circ R(g_2)\qquad\text{for all }g_1,g_2\in G ,
\]
that is, $R$ is a group homomorphism. The map $R$ should also be smooth (the lecture does not say it explicitly [?]). \ts{17}{1}{1:18:08} \ts{17}{1}{1:18:29} \ts{17}{1}{1:18:50} \ts{17}{1}{1:19:10} \ts{17}{1}{1:19:31} \ts{17}{1}{1:19:51} \ts{17}{1}{1:20:12} \ts{17}{1}{1:20:33}
\end{definition}
It looks very similar to the algebra case, but $GL(V)$ is ``a little less'' than $\mathrm{End}(V)$: the requirement $\det\ne0$ is needed so that the images of group elements, which have inverses, are themselves invertible.

\textbf{Example: $SO(2)$ on $\R^2$} \ts{17}{1}{1:20:55}. As a manifold $SO(2)$ is $S^1$, and the group operation moves points along the circle \ts{17}{1}{1:21:17}. To write down a representation one uses a chart $\varphi\colon U\to\R$, the angle map, defined on the circle with one point removed (otherwise the domain would not be open and one cannot cover the circle by one chart) \ts{17}{1}{1:21:37} \ts{17}{1}{1:21:58} \ts{17}{1}{1:22:18}. The group element with coordinate $\phi$ is $\varphi^{-1}(\phi)$, and it is represented by the rotation matrix
\[
  R\bigl(\varphi^{-1}(\phi)\bigr)=\begin{pmatrix}\cos\phi&\sin\phi\\-\sin\phi&\cos\phi\end{pmatrix}.
\]
By the addition theorems for cosine and sine, adding angles equals composing matrices:
\[
  R\bigl(\varphi^{-1}(\phi_1+\phi_2)\bigr)=R\bigl(\varphi^{-1}(\phi_1)\bigr)\circ R\bigl(\varphi^{-1}(\phi_2)\bigr),
\]
so this is a representation of the Lie group $SO(2)$ on $\R^2$, which is usually how one first meets it \ts{17}{1}{1:22:39} \ts{17}{1}{1:23:21} \ts{17}{1}{1:23:42} \ts{17}{1}{1:24:02} \ts{17}{1}{1:24:23} \ts{17}{1}{1:24:44}. The inverse chart is only needed to say which group elements are represented: on groups this is just more complicated than on $\R^n$.
\begin{coursenote}[Not said in class: the whole group, and the group law]
The lecture says one has a representation ``not on the entire group, because one chooses coordinates to write it down, but on $U$''. In fact the matrix is $2\pi$-periodic in $\phi$, so it is well defined also at the removed point and $R$ is defined on all of $SO(2)$. The relation $\phi_1+\phi_2$ on the left is in $\R$ only as long as the sum stays in the chart; globally it is addition mod $2\pi$. \ts{17}{1}{1:22:39}
\end{coursenote}
\textbf{The adjoint representation of a Lie group} \ts{17}{1}{1:25:07}. Are there non-trivial representations of every Lie group? Yes, and we can construct one that goes by the same name as the one for algebras, written with a capital $\mathrm{Ad}$ (for the group), versus the small $\mathrm{ad}$ (for the algebra): they are related, but a priori different. It is built in two steps \ts{17}{1}{1:25:28} \ts{17}{1}{1:25:48}.
\begin{definition}[New concepts: the adjoint map and the adjoint representation of a Lie group]\label{def:rp-Ad}
For $g\in G$ the \textbf{adjoint map} is
\[
  \mathrm{Ad}_g\colon G\to G,\qquad \mathrm{Ad}_g(h):=g\bullet h\bullet g^{-1},
\]
a left translation with $g$ followed by a right translation with $g^{-1}$ (not a linear map: there is no linearity of any kind here). \ts{17}{1}{1:26:10} \ts{17}{1}{1:26:31} \ts{17}{1}{1:26:53} It fixes the identity: $\mathrm{Ad}_g(e)=g\bullet e\bullet g^{-1}=e$ \ts{17}{1}{1:27:13} \ts{17}{1}{1:27:34}. So its push-forward (\cref{def:tb-push}) at the identity goes from the tangent space at $e$ to the tangent space at $\mathrm{Ad}_g(e)=e$ again \ts{17}{1}{1:27:54} \ts{17}{1}{1:28:15} \ts{17}{1}{1:28:36}:
\[
  (\mathrm{Ad}_g)_*\colon T_eG\to T_eG ,
\]
a linear map from the space isomorphic to the Lie algebra (\cref{thm:ll-iso}) to itself \ts{17}{1}{1:28:56}. It is invertible, since its inverse is $(\mathrm{Ad}_{g^{-1}})_*$, because $\mathrm{Ad}_{g^{-1}}$ is the inverse of $\mathrm{Ad}_g$. The \textbf{adjoint representation of the Lie group} is
\[
  \mathrm{Ad}\colon G\to GL(T_eG),\qquad g\mapsto(\mathrm{Ad}_g)_* .
\]
\ts{17}{1}{1:29:17} \ts{17}{1}{1:29:39} \ts{17}{1}{1:30:00} \ts{17}{1}{1:30:21}
\end{definition}
So a group has always a representation on the vector space $T_eG\cong L(G)$ that comes with it: one can find a representation of the Lie algebra on the Lie algebra (\cref{sec:rp-triv}), and a representation of the Lie group on the Lie algebra. If one has no imagination and is just given a Lie group and knows about the Lie algebra, ``with a group I always have a vector space that comes with it'', and the map that does the job is $\mathrm{Ad}$ \ts{17}{1}{1:30:41} \ts{17}{1}{1:31:02} \ts{17}{1}{1:31:22}.
\begin{coursenote}[Not said in class: why it is a representation, and how $\mathrm{Ad}$ and $\mathrm{ad}$ are related]
Homomorphism: $\mathrm{Ad}_{g_1g_2}=\mathrm{Ad}_{g_1}\circ\mathrm{Ad}_{g_2}$, so by the chain rule for push-forwards $(\mathrm{Ad}_{g_1g_2})_*=(\mathrm{Ad}_{g_1})_*\circ(\mathrm{Ad}_{g_2})_*$; and $g\mapsto(\mathrm{Ad}_g)_*$ is smooth because the group operations are. The relation between the two: the derivative of $\mathrm{Ad}$ at the identity is $\mathrm{ad}$, that is, $\mathrm{ad}_a=\frac{\mathrm d}{\mathrm dt}\big|_{t=0}\mathrm{Ad}_{\exp(ta)}$ (``a representation of the group differentiates to a representation of the algebra'', a statement for \cref{ch:l18}; no proof given). \ts{17}{1}{1:25:28}
\end{coursenote}
The lecture closes: finite-dimensional Lie algebras and Lie groups have been discussed; the next time one finally goes back from the algebra up to the group (and, in general, one \emph{cannot} hit the entire group), through the \emph{exponential map}, a concept from differential geometry which on Lie groups has even nicer properties than on general smooth manifolds \ts{17}{1}{1:31:43} \ts{17}{1}{1:32:03} \ts{17}{1}{1:32:23}. This is \cref{ch:l18}.

% ==================================================================
\section*{Key concepts}
\begin{keyconcepts}
\textbf{Retrieval questions.}
\begin{enumerate}
  \item State the definition of a representation of a Lie algebra. Which bracket appears on each side?
  \item Give three different representations of Lie algebras from the lecture, and say on which spaces they live.
  \item When is a representation reducible? Why must the invariant subspace be non-trivial? What is a faithful representation?
  \item Define $K_\rho$ and the Casimir operator $\Omega_\rho$. What do they depend on?
  \item Prove that $c_\rho=\dim L/\dim V$ for an irreducible $\rho$. Where is Schur's lemma used?
  \item Compute $K_\rho$ and $\Omega_\rho$ for $\rho_{\mathrm{spin}}$ of $\mathfrak{so}(3)$. Where does the factor $2$ with respect to physics come from?
  \item Define a representation of a Lie group and the group $GL(V)$. Give the example of $SO(2)$ and the adjoint representation $\mathrm{Ad}$.
\end{enumerate}
\textbf{Concepts.}
\begin{itemize}
  \item Representation of a Lie algebra, representation space (\cref{def:rp-rep}); examples (\cref{sec:rp-examples}); trivial and adjoint representation (\cref{sec:rp-triv}).
  \item Reducible, irreducible, faithful (\cref{def:rp-irred,def:rp-faithful}).
  \item Killing form of a representation, Casimir operator, Schur's lemma, $c_\rho=\dim L/\dim V$ (\cref{def:rp-killing,def:rp-casimir,thm:rp-schur,thm:rp-value}).
  \item Representation of a Lie group and $\mathrm{Ad}$ (\cref{def:rp-grouprep,def:rp-Ad}).
\end{itemize}
\end{keyconcepts}

\flushside
\end{paracol}
